\section{Computational Verification and Historical Application}
\label{sec:evaluation}
We evaluate the algorithms on constructed instances, replay saved forecasts,
replicate the method on an external public source with fresh model fits,
and measure retrieval on a public API. LP checks use reported floating-point
tolerances; counter coefficients and adverse distributions also receive
exact rational checks. The same author conducted both source studies.

\subsection{Comparative formulations and scaling}
The scaling grid has 24 base aggregate instances with one, two or four
blocks, six or twelve coordinates per block, budgets near one quarter or
one half of the coordinates, and two joint coefficient/total profiles.
Four additional one-block profiles use 24 or 48 coordinates, giving 28 in total. The cases
and enumeration rule were fixed before that audit. All 28 compact solves
completed; the 12 prescribed finite-game comparisons agreed within
$1.27\times10^{-15}$ after scaling by $\max(1,W_0)$. The other 16 cases
were skipped by the enumeration rule, not observed to fail.

Table~\ref{tab:risk-solver-comparison} exposes both methods' measured
construction/solve times and whole-worker peak memory on the same eligible
cases, including cases where enumeration is faster. This is comparative
evidence against the enumerated formulation, not against every alternative
compact solver. The larger compact-only runs establish operability within
this grid rather than a measured speedup over an unexecuted competitor.
Detailed per-case dimensions, residuals and outcomes are bundled.

